\scanpage{79}
commutativité dans
\diagram{\begin{tikzcd}[ampersand replacement=\&,column sep=large,row sep=large]
\& \operatorname{Rep}(\Pi_\infty) \arrow[r,no head,"\varphi_{\mathrm{rep}}"] \arrow[d,"\mathrm{pr}_1"] \arrow[dd,bend right=45,"\mathrm{pr}_2"'] \& \operatorname{Drap.aff}(E) \arrow[d,"\mathrm{pr}_1"] \arrow[dd,bend right=40,"\mathrm{pr}_2"] \\
(7) \& S(\Pi_\infty) \arrow[r,"\varphi_s","{\scriptscriptstyle\bullet}"{pos=.79,description,inner sep=0pt,fill=none}] \& E\,\unsure{\rbrace|)} \\
\& A(\Pi_\infty) \arrow[r,no head,"\varphi_a"] \& \operatorname{Dr}(E)
\end{tikzcd}}

La condition [pour $\varphi_G$ \inserted{(cf (1) ou (2))} et le $G$-hom.
$\varphi_{\mathrm{rep}}:\operatorname{Rep}(\Pi_\infty)\unsure{\relbar|\to}\operatorname{Drap.aff}(E)$ donnés, cette
dernière donnée correspondant à la donnée
de (\unsure{4|3})] pour que $\mathrm{pr}_i\circ\varphi_{\mathrm{rep}}$ passe au quotient \inserted{en $\varphi_s$ et $\varphi_a$}
-- ou encore pour que $\varphi_s$ et $\varphi_a$ soient
compatibles avec les relations d'incidence,
c'est que l'on ait
\[ \sigma_1\,\mathrm{pr}_1(r)=\mathrm{pr}_1(r),\quad \sigma_0\,\mathrm{pr}_2(r)=\mathrm{pr}_{\unsure{2|1}}(r) \]
i.e.
\[ \boxed{\sigma_1s_0=s_0,\ \sigma_0a_0=a_0} \tag*{\deleted{\unsure{(7)}}(8)} \]

Il reste à exprimer que les $s_n$ engendrent
affinement $E$. On va introduire
% archive sheet number: 78
